\chapter{Meerelektronatomen en termsymbolen}\label{ch:b3:many-electron-atoms}

In 1868 onthulde een gele lijn in het spectrum van de zon een element dat
op aarde onbekend was: helium. Toen het in 1895 eindelijk in het
laboratorium werd geïsoleerd, leek zijn spectrum op dat van \emph{twee}
elementen: twee families van
lijnen, die een tijdlang werden toegeschreven aan ``orthohelium'' en
``parahelium'', en die nooit licht met elkaar uitwisselden. Er bestaat maar
één helium. De twee families zijn de triplet- en \omterm{def:b3:many-electron-atoms:singlet-triplet}{singulettoestanden} van
hetzelfde atoom, gescheiden gehouden door de spin van het elektron en door
een regel die fundamenteler is dan welke kracht ook: de golffunctie van
twee elektronen moet van teken veranderen wanneer ze verwisseld worden.
Dit hoofdstuk volgt die regel van het uitsluitingsprincipe van Pauli tot
de termsymbolen waarmee spectroscopisten elke toestand van elk atoom en
ion benoemen --- ook die van de overgangsmetaalionen waarvan de kleuren
verklaard worden
in \cref{ch:b3:complex-spectra-magnetism}.

\begin{recall}
Het deel van jaar~1 kende elektronen vier kwantumgetallen toe, $n$, $l$,
$m_l$, $m_s$, bouwde grondtoestandsconfiguraties op met het
uitsluitingsprincipe van Pauli, de regel van Klechkowski en de regel van
Hund, en telde ongepaarde elektronen. Het deel van jaar~2
gaf orbitaalenergieën en de regels van Slater voor de afscherming.
\Cref{ch:b3:quantum-model-systems} leverde \omterm{def:b3:quantum-model-systems:operator}{operatoren}, \omterm{def:b3:quantum-model-systems:operator}{eigenwaarden}, het
impulsmoment van de rotor ($\hat L^2$ met \omterm{def:b3:quantum-model-systems:operator}{eigenwaarden} $l(l+1)\hbar^2$) en het
waterstofatoom.
\end{recall}

\section{Spin en spinorbitalen}

Het elektron draagt een intrinsiek impulsmoment, zijn spin, met
kwantumgetal $s = \frac12$: $\hat S^2$ heeft de enige \omterm{def:b3:quantum-model-systems:operator}{eigenwaarde} $s(s+1)\hbar^2 =
\frac34\hbar^2$, en $\hat S_z$ de twee \omterm{def:b3:quantum-model-systems:operator}{eigenwaarden} $m_s\hbar = \pm\frac12\hbar$.
De twee spinfuncties worden geschreven als $\alpha$ ($m_s = +\frac12$) en $\beta$
($m_s = -\frac12$); ze zijn orthonormaal, $\langle\alpha|\alpha\rangle =
\langle\beta|\beta\rangle = 1$ en $\langle\alpha|\beta\rangle = 0$. Spin heeft
geen klassieke tegenhanger; het is niet de rotatie van een bolletje, en het
komt in een niet-relativistische \omterm{def:b3:quantum-model-systems:schrodinger}{hamiltonoperator} helemaal niet voor.

\begin{definition}[Spinorbitaal]\label{def:b3:many-electron-atoms:spin-orbital}
Een \emph{spinorbitaal}\index{spinorbitaal} is het product van een
ruimtelijke orbitaal
$\phi(\mathbf r)$ en een spinfunctie: $\phi\alpha$ of $\phi\beta$. Een
spinorbitaal bevat de volledige beschrijving van één elektron.
\end{definition}

Een atoomorbitaal die in het deel van jaar~1 met $(n, l, m_l)$ werd
aangeduid, levert dus twee
spinorbitalen; daarom bevat een subschil $2(2l+1)$ elektronen. De echte
vraag is waarom een \omterm{def:b3:many-electron-atoms:spin-orbital}{spinorbitaal} er hoogstens één bevat.

\section{Ononderscheidbare elektronen en de slaterdeterminant}

\begin{definition}[Ononderscheidbare deeltjes, antisymmetrische golffunctie]\label{def:b3:many-electron-atoms:indistinguishable}
Deeltjes zijn \emph{ononderscheidbaar}\index{ononderscheidbare deeltjes} als
geen enkele meting ze uit elkaar kan houden: elk elektron is identiek aan
elk ander. Een golffunctie van meerdere elektronen is een
\emph{antisymmetrische golffunctie}\index{antisymmetrische golffunctie} als
ze van teken verandert wanneer de
coördinaten (ruimte en spin) van twee willekeurige elektronen verwisseld
worden:
$\Psi(\dots,i,\dots,j,\dots) = -\Psi(\dots,j,\dots,i,\dots)$.
\end{definition}

Ononderscheidbaarheid alleen eist dat $|\Psi|^2$ onveranderd blijft bij
een verwisseling, zodat $\Psi$ hoogstens van teken kan veranderen. De
natuur heeft gekozen: de
golffunctie van elk systeem van elektronen --- van alle fermionen --- is
antisymmetrisch. Dit is een postulaat van de kwantummechanica, bevestigd
door elk atoomspectrum.

\begin{definition}[Slaterdeterminant]\label{def:b3:many-electron-atoms:slater}
Voor $N$ elektronen in $N$ verschillende spinorbitalen $\chi_1, \dots, \chi_N$ is de
\emph{slaterdeterminant}\index{slaterdeterminant}
\[
\Psi = \frac{1}{\sqrt{N!}}\begin{vmatrix}
\chi_1(1) & \chi_2(1) & \cdots & \chi_N(1)\\
\chi_1(2) & \chi_2(2) & \cdots & \chi_N(2)\\
\vdots & & & \vdots\\
\chi_1(N) & \chi_2(N) & \cdots & \chi_N(N)
\end{vmatrix},
\]
waarin rij $i$ elektron $i$ in elk \omterm{def:b3:many-electron-atoms:spin-orbital}{spinorbitaal} bevat. Kort geschreven:
$|\chi_1\chi_2\dots\chi_N|$.
\end{definition}

\begin{theorem}[Het uitsluitingsprincipe van Pauli]\label{thm:b3:many-electron-atoms:pauli}
Een \omterm{def:b3:many-electron-atoms:slater}{slaterdeterminant} is antisymmetrisch, en hij is nul als twee
elektronen hetzelfde \omterm{def:b3:many-electron-atoms:spin-orbital}{spinorbitaal} bezetten. Bijgevolg kunnen geen twee
elektronen van een atoom dezelfde vier kwantumgetallen hebben.
\end{theorem}

\begin{proof}
Het verwisselen van de elektronen $i$ en $j$ verwisselt twee rijen van de
determinant, wat zijn teken omkeert: de antisymmetrie geldt voor elk paar.
Als twee
spinorbitalen gelijk zijn, zijn twee kolommen gelijk en is de determinant
nul: zo'n toestand bestaat niet. Twee elektronen van hetzelfde atoom met
dezelfde $n$, $l$,
$m_l$ en $m_s$ zouden in hetzelfde \omterm{def:b3:many-electron-atoms:spin-orbital}{spinorbitaal} zitten.
\end{proof}

\begin{example}[De grondtoestand van helium]\label{ex:b3:many-electron-atoms:helium-ground}
Met beide elektronen in $1s$ is
\[
\Psi = \frac{1}{\sqrt2}\begin{vmatrix}1s\alpha(1) & 1s\beta(1)\\ 1s\alpha(2) &
1s\beta(2)\end{vmatrix} = 1s(1)\,1s(2)\cdot\frac{1}{\sqrt2}\big[\alpha(1)\beta(2)
- \beta(1)\alpha(2)\big].
\]
Het ruimtelijke deel is symmetrisch, het spindeel antisymmetrisch: de twee
spins zijn gepaard, met totale spin nul.
\end{example}

\section{Helium: coulombafstoting en uitwisseling}

Breng één elektron van helium naar $2s$. Met $1s\alpha$, $1s\beta$,
$2s\alpha$, $2s\beta$ kunnen vier determinanten gebouwd worden met één
elektron in elke orbitaal.
Herschikt worden het producten van één ruimtelijke en één spinfunctie:
\[
\Psi_\pm = \frac{1}{\sqrt2}\big[1s(1)2s(2) \pm 2s(1)1s(2)\big]\times(\text{spinfunctie}).
\]
De antisymmetrische ruimtelijke functie $\Psi_-$ moet gecombineerd worden
met een symmetrische spinfunctie, en daarvan zijn er drie: $\alpha(1)\alpha(2)$,
$\beta(1)\beta(2)$ en $\frac{1}{\sqrt2}[\alpha(1)\beta(2) + \beta(1)\alpha(2)]$, de
drie componenten $M_S = 1, 0, -1$ van een totale spin $S = 1$. De
symmetrische $\Psi_+$ moet gecombineerd worden met de enige
antisymmetrische spinfunctie, $S = 0$.

\begin{definition}[Singulet- en triplettoestanden]\label{def:b3:many-electron-atoms:singlet-triplet}
Een toestand met totale spin $S = 0$ is een
\emph{singulettoestand}\index{singulettoestand}; een toestand met totale
spin $S = 1$, waarvan de drie componenten $M_S = 1, 0, -1$ zonder
magneetveld dezelfde energie hebben, is een
\emph{triplettoestand}\index{triplettoestand}.
\end{definition}

\begin{definition}[Uitwisselingsintegraal]\label{def:b3:many-electron-atoms:exchange}
Voor twee orbitalen $a$ en $b$ en de elektronenafstoting $e^2/4\pi\varepsilon_0
r_{12}$ is de \emph{uitwisselingsintegraal}\index{uitwisselingsintegraal}
\[
K_{ab} = \iint a(1)b(2)\,\frac{e^2}{4\pi\varepsilon_0r_{12}}\,b(1)a(2)\,\dd\tau_1\dd\tau_2,
\]
dezelfde integraal als de klassieke elektronenafstoting $J_{ab}$ (met
$a(1)b(2)$ aan beide kanten), behalve dat de elektronen aan één kant
verwisseld zijn. $K_{ab}$ is positief en heeft geen klassieke tegenhanger.
\end{definition}

\begin{proposition}[Energieën van singulet en triplet]\label{prop:b3:many-electron-atoms:helium}
In eerste orde in de elektronenafstoting hebben de toestanden $1s2s$ van
helium de energieën $E_\pm = E_{1s} + E_{2s} + J \pm K$; de singulet ($+$)
ligt boven de
triplet ($-$), met een verschil van $2K$.
\end{proposition}

\begin{proof}
De ongestoorde energie van beide $\Psi_\pm$ is $E_{1s} + E_{2s}$
(waterstofachtige orbitalen met kernlading 2). De eersteordecorrectie is
de \omterm{def:b3:quantum-model-systems:hermitian}{verwachtingswaarde} van de afstoting $\hat V_{12}$ in de genormeerde toestand:
$\frac12\langle 1s(1)2s(2) \pm 2s(1)1s(2)|\hat V_{12}|1s(1)2s(2) \pm
2s(1)1s(2)\rangle$. De twee directe termen geven $\frac12(J + J)$; de twee
kruistermen geven $\pm\frac12(K + K)$, omdat $\hat V_{12}$ symmetrisch is in de
elektronen. Dus $\Delta E = J \pm K$, en de spinfuncties, genormeerd en
niet beïnvloed door $\hat V_{12}$, vermenigvuldigen alleen met 1.
\end{proof}

De triplet ligt lager, hoewel de \omterm{def:b3:quantum-model-systems:schrodinger}{hamiltonoperator} geen spin bevat: zijn
ruimtelijke functie is nul wanneer $\mathbf r_1 = \mathbf r_2$, zodat de
twee elektronen elkaar ontwijken en elkaar minder afstoten. Dit
``fermigat'' is de fysische inhoud van de eerste regel van Hund.

\begin{example}[De uitwisselingsintegraal van helium, gemeten]\label{ex:b3:many-electron-atoms:K}
% ledger: lev:He.2s3S1, lev:He.2s1S0
De niveaus $1s2s$ van helium liggen bij \qty{159856}{cm^{-1}} (${}^3$S) en
\qty{166277}{cm^{-1}} (${}^1$S) boven de grondtoestand: de singulet ligt
\qty{6421}{cm^{-1}}, of \qty{0.796}{eV}, hoger, en $K(1s,2s) = \qty{0.398}{eV}$.
De figuur verderop toont de twee families.
\end{example}

\section{Russell-saunderskoppeling en termsymbolen}

In een licht atoom tellen de baanimpulsmomenten van de elektronen op tot
een totaal $\mathbf L$, hun spins tot een totaal $\mathbf S$, en pas daarna
koppelen
$\mathbf L$ en $\mathbf S$ zwak met elkaar.

\begin{definition}[Russell-saunderskoppeling]\label{def:b3:many-electron-atoms:russell-saunders}
Bij \emph{russell-saunderskoppeling}\index{russell-saunderskoppeling} worden de
toestanden van een configuratie aangeduid met het \emph{kwantumgetal van
het totale baanimpulsmoment}\index{kwantumgetal van het totale baanimpulsmoment}
$L$ ($\hat{\mathbf L}^2 = L(L+1)\hbar^2$, met $M_L = \sum m_l$), het \emph{totale
spinkwantumgetal}\index{totale spinkwantumgetal} $S$ ($M_S = \sum m_s$),
en het \emph{kwantumgetal van het totale impulsmoment}\index{kwantumgetal van het totale impulsmoment}
$J$, dat de waarden $|L - S|, \dots, L + S$ aanneemt.
\end{definition}

\begin{definition}[Spectroscopische term, spinmultipliciteit, termsymbool]\label{def:b3:many-electron-atoms:term}
Een \emph{spectroscopische term}\index{spectroscopische term} is de
verzameling van de
$(2L+1)(2S+1)$ toestanden van een configuratie met gegeven $L$ en $S$. Zijn
\emph{spinmultipliciteit}\index{spinmultipliciteit} is $2S + 1$. Het
\emph{termsymbool}
\index{termsymbool} is ${}^{2S+1}L$, met de letters S, P, D, F, G, H
voor $L = 0, 1, 2, 3, 4, 5$; een toestand met gegeven $J$ wordt geschreven als ${}^{2S+1}L_J$ ---
bijvoorbeeld \termsym{3}{P}{2}.
\end{definition}

\begin{proposition}[Gesloten schillen]\label{prop:b3:many-electron-atoms:closed-shell}
Een gevulde subschil heeft $L = 0$ en $S = 0$; ze draagt niets bij aan het
\omterm{def:b3:many-electron-atoms:term}{termsymbool}.
\end{proposition}

\begin{proof}
In een gevulde subschil komt elke $m_l$ tweemaal voor en elke $m_s = \pm\frac12$
$2l+1$ keer, zodat $M_L = 0$ en $M_S = 0$; en er is maar één manier om ze
te vullen (één determinant). Eén enkele toestand met $M_L = M_S = 0$ kan alleen
behoren tot $L = 0$, $S = 0$.
\end{proof}

\begin{proposition}[De toestanden van een configuratie tellen]\label{prop:b3:many-electron-atoms:count}
$N$ elektronen in een subschil van $2(2l+1)$ spinorbitalen geven
$\binom{2(2l+1)}{N}$ determinanten, en de ontaardingsgraden $(2L+1)(2S+1)$ van
haar termen tellen op tot dit getal.
\end{proposition}

\begin{proof}
Een determinant ligt vast door de verzameling bezette spinorbitalen (hun
volgorde verandert alleen zijn teken): men kiest er $N$ uit de $2(2l+1)$.
De termen zijn dezelfde toestanden, anders gegroepeerd, dus de aantallen
moeten overeenstemmen.
\end{proof}

\begin{method}[De termen van een configuratie]\label{met:b3:many-electron-atoms:terms}
\begin{enumerate}
  \item Som de determinanten op die het uitsluitingsprincipe toelaat en
        rangschik ze in een tabel naar $M_L$ en $M_S$.
  \item Zoek de grootste $M_L$; met de grootste $M_S$ die erbij hoort,
        begint die een term met $L = M_L$, $S = M_S$.
  \item Schrap één determinant voor elk paar $(M_L, M_S)$ met
        $|M_L| \le L$, $|M_S| \le S$.
  \item Herhaal met wat overblijft tot de tabel leeg is; controleer het aantal.
\end{enumerate}
\end{method}

\begin{example}[De termen van $p^2$]\label{ex:b3:many-electron-atoms:p2}
Twee elektronen in $p$ ($m_l = 1, 0, -1$) geven $\binom62 = 15$ determinanten.
Hun tabel naar $M_L$ en $M_S$ is:
\begin{center}\small
\begin{tabular}{c|ccc}
\hline
 & $M_S = 1$ & $M_S = 0$ & $M_S = -1$\\ \hline
$M_L = 2$ & -- & $(1^+,1^-)$ & --\\
$M_L = 1$ & $(1^+,0^+)$ & $(1^+,0^-)$, $(1^-,0^+)$ & $(1^-,0^-)$\\
$M_L = 0$ & $(1^+,-1^+)$ & $(1^+,-1^-)$, $(1^-,-1^+)$, $(0^+,0^-)$ & $(1^-,-1^-)$\\
$M_L = -1$ & $(0^+,-1^+)$ & $(0^+,-1^-)$, $(0^-,-1^+)$ & $(0^-,-1^-)$\\
$M_L = -2$ & -- & $(-1^+,-1^-)$ & --\\ \hline
\end{tabular}
\end{center}
(met $m_l^\pm$ voor $m_s = \pm\frac12$). $M_L = 2$ komt alleen voor met $M_S = 0$:
een term ${}^1$D (5 toestanden). De grootste resterende $M_L = 1$ komt voor met
$M_S = 1$: een term ${}^3$P (9 toestanden). Er blijft één toestand over met $M_L = M_S = 0$: een
term ${}^1$S. Inderdaad is $5 + 9 + 1 = 15$.
\end{example}

\begin{example}[Koolstof, gemeten]\label{ex:b3:many-electron-atoms:carbon}
% ledger: lev:C.3P1, lev:C.3P2, lev:C.1D2, lev:C.1S0
De grondtoestandsconfiguratie $2p^2$ van koolstof geeft de niveaus \termsym{3}{P}{0},
\termsym{3}{P}{1}, \termsym{3}{P}{2} bij 0, 16.4 en \qty{43.4}{cm^{-1}}, dan
\termsym{1}{D}{2} bij \qty{10193}{cm^{-1}} en \termsym{1}{S}{0} bij
\qty{21648}{cm^{-1}}: de term ${}^3$P ligt het laagst, zoals de regels van
Hund voorspellen.
\end{example}

\begin{omfigure}
\begin{tikzpicture}
\begin{axis}[name=main, width=0.62\linewidth, height=6.4cm, xmin=-0.3, xmax=6.6,
  ymin=-1500, ymax=23500, axis y line=left, axis x line=none,
  ylabel={energie / cm$^{-1}$}, unbounded coords=jump, clip=false,
  scaled y ticks=false, ytick={0,10000,20000},
  tick label style={font=\small}, label style={font=\small}]
\addplot[omstate] table[x=x, y=E] {figdata/out/bachelor-3-atomic-levels-carbon.dat};
\draw[densely dotted, gray] (axis cs:1,4858.5) -- (axis cs:2,29.6);
\draw[densely dotted, gray] (axis cs:1,4858.5) -- (axis cs:2,10192.7);
\draw[densely dotted, gray] (axis cs:1,4858.5) -- (axis cs:2,21648);
\draw[densely dotted, gray] (axis cs:3,10192.7) -- (axis cs:4,10192.7);
\draw[densely dotted, gray] (axis cs:3,21648) -- (axis cs:4,21648);
\node[font=\small, anchor=south] at (axis cs:0.5,4858.5) {$2p^2$};
\node[font=\small, anchor=south] at (axis cs:2.5,21648) {${}^1$S};
\node[font=\small, anchor=south] at (axis cs:2.5,10192.7) {${}^1$D};
\node[font=\small, anchor=south] at (axis cs:2.5,29.6) {${}^3$P};
\node[font=\small, anchor=west] at (axis cs:5.05,21648) {\termsym{1}{S}{0}};
\node[font=\small, anchor=west] at (axis cs:5.05,10192.7) {\termsym{1}{D}{2}};
\node[font=\scriptsize, anchor=north] at (axis cs:0.5,-600) {configuratie};
\node[font=\scriptsize, anchor=north] at (axis cs:2.5,-600) {termen};
\node[font=\scriptsize, anchor=north] at (axis cs:4.5,-600) {niveaus};
\draw[gray, ->] (axis cs:3.05,60) -- (axis cs:6.9,2500);
\end{axis}
\begin{axis}[at={(main.south east)}, anchor=south west, xshift=1.2cm, yshift=0.9cm,
  width=3.6cm, height=3.6cm, xmin=-0.1, xmax=1.1, ymin=-5, ymax=50,
  axis y line=right, axis x line=none, unbounded coords=jump, ytick={0,16.4,43.4},
  yticklabels={0,16.4,43.4}, tick label style={font=\scriptsize}, clip=false,
  title={\scriptsize ${}^3$P, vergroot}, title style={yshift=-4pt}]
\addplot[omstate] table[x=x, y=E] {figdata/out/bachelor-3-atomic-levels-carbon-zoom.dat};
\node[font=\scriptsize, anchor=east] at (axis cs:0,0) {$J = 0$};
\node[font=\scriptsize, anchor=east] at (axis cs:0,16.4) {$J = 1$};
\node[font=\scriptsize, anchor=east] at (axis cs:0,43.4) {$J = 2$};
\end{axis}
\end{tikzpicture}

{\small De grondtoestandsconfiguratie van koolstof: één configuratie, drie
termen (bij hun gemeten energieën; de term ${}^3$P bij het gewogen
gemiddelde van zijn niveaus, de configuratie bij het gemiddelde van alle
vijftien toestanden) en vijf niveaus. De spin-baanopsplitsing van ${}^3$P
(vergroot, in
\unit{cm^{-1}}) is enkele honderden keren kleiner dan de afstanden tussen
de termen.}
\end{omfigure}

\section{De regels van Hund en spin-baankoppeling}

\begin{proposition}[Regels van Hund voor termen]\label{prop:b3:many-electron-atoms:hund-terms}
Voor de grondtoestandsconfiguratie van een atoom of ion is de laagste term
die met (1) de grootste $S$, en (2) bij die $S$ de grootste $L$. Het laagste
niveau ervan heeft (3) $J = |L - S|$ als de subschil minder dan half
gevuld is, en $J = L +
S$ als ze meer dan half gevuld is.
\end{proposition}
\begin{proof}[Status]
Deze regels zijn experimenteel vastgesteld, en alleen voor
grondtoestandsconfiguraties; het zijn geen stellingen. Regel 1 is de
stabilisatie door uitwisseling uit de vorige paragraaf; regel 2 zegt dat
elektronen die in dezelfde zin draaien elkaar minder vaak ontmoeten;
regel 3 volgt uit het teken van de \omterm{def:b3:many-electron-atoms:spin-orbit}{spin-baankoppeling}, hieronder.
\end{proof}

\begin{method}[De grondterm uit een hokjesschema]\label{met:b3:many-electron-atoms:ground-term}
\begin{enumerate}
  \item Vul de hokjes van de open subschil vanaf $m_l = +l$ naar beneden,
        één elektron per hokje met parallelle spins, en paar daarna
        opnieuw vanaf $+l$.
  \item $S = \frac12 \times$(aantal ongepaarde elektronen); $L = |\sum m_l|$.
  \item $J = |L - S|$ bij minder dan halfvol, $L + S$ bij meer dan halfvol;
        bij precies halfvol is $L = 0$ en $J = S$.
\end{enumerate}
\end{method}

\begin{example}[Grondtermen]\label{ex:b3:many-electron-atoms:ground-terms}
% ledger: lev:Ti2.3P0, lev:Ni2.3P0, lev:Cr3.4P1_2
\ce{Ti^2+}, $d^2$: hokjes $m_l = 2, 1$ gevuld, $S = 1$, $L = 3$, $J = 2$:
\termsym{3}{F}{2}. \ce{Cr^3+}, $d^3$: $S = \frac32$, $L = 2 + 1 + 0 = 3$,
\termsym{4}{F}{3/2}. \ce{Fe^2+}, $d^6$: vijf omhoog, één omlaag in $m_l = 2$:
$S = 2$, $L = 2$, meer dan halfvol, \termsym{5}{D}{4}. \ce{Ni^2+}, $d^8$:
$S = 1$, $L = 3$, \termsym{3}{F}{4}. Elk komt overeen met het grondniveau
dat de atoomspectroscopie meet.
\end{example}

\begin{definition}[Spin-baankoppeling, fijnstructuur]\label{def:b3:many-electron-atoms:spin-orbit}
\emph{Spin-baankoppeling}\index{spin-baankoppeling} is de wisselwerking
tussen de magnetische momenten van de spin van een elektron en van zijn
baanbeweging, voor een term geschreven als $A\,\hat{\mathbf L}\cdot\hat{\mathbf S}/\hbar^2$. Ze
splitst een term op in zijn niveaus met verschillende $J$: de
\emph{fijnstructuur}\index{fijnstructuur} van de term.
\end{definition}

\begin{proposition}[Intervalregel van Landé]\label{prop:b3:many-electron-atoms:lande}
Binnen een term is $E(J) = E_0 + \frac{A}{2}[J(J+1) - L(L+1) - S(S+1)]$, zodat
$E(J) - E(J - 1) = AJ$. $A > 0$ (normaal) voor een subschil die minder dan
half gevuld is,
$A < 0$ (omgekeerd) voor een die meer dan half gevuld is.
\end{proposition}

\begin{proof}
$\hat{\mathbf J} = \hat{\mathbf L} + \hat{\mathbf S}$ geeft $\hat{\mathbf J}^2 =
\hat{\mathbf L}^2 + \hat{\mathbf S}^2 + 2\hat{\mathbf L}\cdot\hat{\mathbf S}$, dus
$\hat{\mathbf L}\cdot\hat{\mathbf S} = \frac12(\hat{\mathbf J}^2 - \hat{\mathbf L}^2 -
\hat{\mathbf S}^2)$, met als \omterm{def:b3:quantum-model-systems:operator}{eigenwaarde} in een toestand met gegeven $J$, $L$, $S$
$\frac{\hbar^2}{2}[J(J+1) - L(L+1) - S(S+1)]$. Het verschil tussen $J$ en
$J - 1$ is $\frac A2[J(J+1) - (J-1)J] = AJ$. Het teken van $A$ wordt
aangenomen: een
subschil die meer dan half gevuld is, gedraagt zich als het overeenkomende
aantal positieve gaten.
\end{proof}

\begin{example}[De intervalregel getoetst]\label{ex:b3:many-electron-atoms:lande-test}
% ledger: lev:C.3P1, lev:C.3P2, lev:O.3P1, lev:O.3P0
Voor ${}^3$P voorspelt de regel intervallen in de verhouding $J = 2 : 1$, dus 2.
Koolstof: $(43.4 - 16.4)/16.4 = 1.65$. Zuurstof ($2p^4$, omgekeerd, \termsym{3}{P}{2}
het laagst, dan $J = 1$ bij \qty{158.3}{cm^{-1}} en $J = 0$ bij \qty{227.0}{cm^{-1}}):
$158.3/68.7 = 2.30$. De regel klopt binnen 20\,\%: \omterm{def:b3:many-electron-atoms:russell-saunders}{russell-saunderskoppeling}
is een goede beschrijving van lichte atomen, geen exacte.
\end{example}

\begin{omfigure}
\begin{tikzpicture}[font=\small]
  \draw[omstate] (0,0) -- (1.6,0) node[midway, below] {$3s\ {}^2$S$_{1/2}$};
  \draw[omstate] (0,3.0) -- (1.6,3.0);
  \node[anchor=east] at (-0.1,3.0) {$3p$};
  \draw[omstate] (2.8,2.75) -- (4.4,2.75) node[right] {${}^2$P$_{1/2}$};
  \draw[omstate] (2.8,3.25) -- (4.4,3.25) node[right] {${}^2$P$_{3/2}$};
  \draw[densely dotted, gray] (1.6,3.0) -- (2.8,2.75) (1.6,3.0) -- (2.8,3.25);
  \draw[omfluo] (3.3,2.73) -- (3.3,0.02) node[pos=0.55, left] {D$_1$};
  \draw[omfluo] (3.9,3.23) -- (3.9,0.02) node[pos=0.55, right] {D$_2$};
  \draw[omstate] (2.8,0) -- (4.4,0);
  \draw[densely dotted, gray] (1.6,0) -- (2.8,0);
  \draw[decorate, decoration={brace, amplitude=3pt}] (5.5,3.27) -- (5.5,2.73)
     node[midway, right=3pt, align=left] {\footnotesize \qty{17.2}{cm^{-1}}\\[-1pt]\footnotesize (niet op schaal)};
\end{tikzpicture}

{\small De D-lijnen van natrium. Het $3p$-niveau wordt door
\omterm{def:b3:many-electron-atoms:spin-orbit}{spin-baankoppeling} opgesplitst
in ${}^2$P$_{1/2}$ en ${}^2$P$_{3/2}$; emissie naar het grondniveau $3s$
geeft twee gele lijnen, D$_1$ bij \qty{589.76}{nm} en D$_2$ bij
\qty{589.16}{nm} (golflengten in vacuüm).}
\end{omfigure}

\begin{example}[Het doublet van natrium]\label{ex:b3:many-electron-atoms:sodium}
% ledger: lev:Na.3p1_2, lev:Na.3p3_2
De $3p$-niveaus van natrium liggen bij 16956.17 en \qty{16973.37}{cm^{-1}}: de
D-lijnen liggen bij $10^7/16956.17 = \qty{589.76}{nm}$ en \qty{589.16}{nm} in
vacuüm, op een afstand van \qty{17.20}{cm^{-1}}. Voor ${}^2$P is $E(\frac32) - E(\frac12) =
\frac32A$, dus $A = \qty{11.47}{cm^{-1}}$.
\end{example}

\begin{proposition}[Selectieregels voor atomen]\label{prop:b3:many-electron-atoms:selection-atoms}
Bij \omterm{def:b3:many-electron-atoms:russell-saunders}{russell-saunderskoppeling} vereist een elektrische-dipoolovergang
$\Delta S = 0$, $\Delta L = 0, \pm1$, $\Delta J = 0, \pm1$ (maar niet $J = 0 \to J
= 0$), en een verandering van pariteit: voor de sprong van één elektron $\Delta l = \pm1$.
\end{proposition}
\begin{proof}[Status]
$\Delta S = 0$ wordt bewezen in \cref{ch:b3:electronic-spectroscopy}: de
dipooloperator werkt niet in op de spin. De andere volgen uit het
vectorkarakter van de dipooloperator en worden hier aangenomen.
\end{proof}

\begin{omfigure}
\begin{tikzpicture}
\begin{axis}[width=0.62\linewidth, height=6.2cm, xmin=-0.5, xmax=3.5,
  ymin=158, ymax=173, axis y line=left, axis x line=none,
  ylabel={energie / $10^3$ cm$^{-1}$}, unbounded coords=jump, clip=false,
  tick label style={font=\small}, label style={font=\small}]
\addplot[omstate] table[x=x, y=E] {figdata/out/bachelor-3-atomic-levels-helium.dat};
\node[font=\small, anchor=west] at (axis cs:1.05,166.277) {$1s2s\ {}^1$S};
\node[font=\small, anchor=west] at (axis cs:1.05,171.135) {$1s2p\ {}^1$P};
\node[font=\small, anchor=west] at (axis cs:3.05,159.856) {$1s2s\ {}^3$S};
\node[font=\small, anchor=west] at (axis cs:3.05,169.087) {$1s2p\ {}^3$P};
\node[font=\small, anchor=south] at (axis cs:0.5,172.3) {singuletten};
\node[font=\small, anchor=south] at (axis cs:2.5,172.3) {tripletten};
\draw[omfluo] (axis cs:0.5,171.0) -- (axis cs:0.5,166.42) node[pos=0.5, left, font=\scriptsize] {\qty{2059}{nm}};
\draw[omfluo] (axis cs:2.5,168.95) -- (axis cs:2.5,160.0) node[pos=0.5, left, font=\scriptsize] {\qty{1083}{nm}};
\draw[densely dotted, gray] (axis cs:-0.3,159.856) -- (axis cs:2,159.856);
\draw[<->, omThm] (axis cs:-0.2,159.856) -- (axis cs:-0.2,166.277) node[pos=0.5, right, font=\scriptsize] {$2K$};
\end{axis}
\end{tikzpicture}

{\small De twee ``heliums''. Singulet- en tripletniveaus van de configuraties
$1s2s$ en $1s2p$ bij hun gemeten energieën boven de grondtoestand (de
term ${}^3$P bij het gemiddelde van zijn niveaus). Lijnen verbinden alleen
toestanden van dezelfde familie; elke triplet ligt $2K$ onder zijn
singulet.}
\end{omfigure}

\begin{history}[Twee heliums en het uitsluitingsprincipe]
Helium werd in 1868 naar de zon genoemd en in 1895 door William Ramsay op
aarde gevonden. Zijn twee lijnenreeksen bleven spectroscopisten dertig
jaar lang voor een raadsel stellen. In 1925 stelde Wolfgang Pauli voor dat
geen twee elektronen dezelfde kwantumgetallen delen, in hetzelfde jaar
waarin George Uhlenbeck en Samuel Goudsmit de elektronspin voorstelden; in
1926 toonde Werner Heisenberg aan dat de
symmetrie van de golffunctie alleen al helium in zijn singulet- en
tripletfamilie opsplitst.
\end{history}

\begin{inthelab}[Het natriumdoublet scheiden]
Een natriumlamp wordt voor de spleet van een traliespectrometer
geplaatst. Met een tralie van 600 lijnen per millimeter verschillen de
hoeken van de twee lijnen in de eerste orde ongeveer \num{0.02}\textdegree:
een goede spectrometer scheidt ze, een eenvoudig practicumprisma niet. De
lamp wordt heet; ze wordt pas aangeraakt nadat ze is afgekoeld.
\end{inthelab}

\section{Oefeningen}

\begin{exercise}[$\star$]\label{exo:b3:many-electron-atoms:1}
Schrijf de \omterm{def:b3:many-electron-atoms:slater}{slaterdeterminant} op van de grondtoestand van lithium,
$1s^22s$, met het elektron in $2s$ in spin $\alpha$. Waarom bestaat er
geen determinant voor
$1s^3$?
\end{exercise}

\begin{exercise}[$\star$]\label{exo:b3:many-electron-atoms:2}
Tel de determinanten van de $p^2$-configuraties, $p^3$ en $d^2$. De termen
van $p^3$ zijn ${}^4$S, ${}^2$D, ${}^2$P en die van $d^2$ zijn ${}^3$F, ${}^3$P, ${}^1$G,
${}^1$D, ${}^1$S: controleer beide aantallen.
\end{exercise}

\begin{exercise}[$\star$]\label{exo:b3:many-electron-atoms:3}
Geef de grondterm en het grondniveau van N, O, F, \ce{Ti^2+} en \ce{Fe^3+} ($d^5$).
\end{exercise}

\begin{exercise}[$\star$]\label{exo:b3:many-electron-atoms:4}
Geef de niveaus van de termen ${}^2$D en ${}^3$P met hun ontaardingsgraden $2J+1$,
en controleer dat de ontaardingsgraden optellen tot $(2L+1)(2S+1)$.
\end{exercise}

\begin{exercise}[$\star\star$]\label{exo:b3:many-electron-atoms:5}
De niveaus ${}^3$P van koolstof liggen bij 0, 16.4 en \qty{43.4}{cm^{-1}}. Bereken de
verhouding van de intervallen en de spin-baankoppelingsconstante $A$ uit
elk interval. Wat zegt het verschil tussen de twee waarden van $A$?
\end{exercise}

\begin{exercise}[$\star\star$]\label{exo:b3:many-electron-atoms:6}
De $3p$-niveaus van natrium liggen bij 16956.17 en \qty{16973.37}{cm^{-1}}. Bereken
de golflengten in vacuüm van de D-lijnen, hun afstand in \unit{cm^{-1}} en
in meV, en $A$ voor de term $3p$.
\end{exercise}

\begin{exercise}[$\star\star$]\label{exo:b3:many-electron-atoms:7}
Welke van deze overgangen van natrium zijn toegestaan in emissie: $3p \to 3s$,
$3d \to 3s$, $3d \to 3p$, $4s \to 3s$, $4s \to 3p$? Verantwoord elk antwoord.
\end{exercise}

\begin{exercise}[$\star\star$]\label{exo:b3:many-electron-atoms:8}
% ledger: lev:He.2p3P2, lev:He.2p3P1, lev:He.2p3P0, lev:He.2p1P1
De niveaus $1s2p$ van helium zijn \termsym{3}{P}{2}, \termsym{3}{P}{1},
\termsym{3}{P}{0} bij 169086.76, 169086.84, \qty{169087.83}{cm^{-1}} en
\termsym{1}{P}{1} bij \qty{171134.89}{cm^{-1}}. Bereken de gemiddelde energie
van de term ${}^3$P, en daarna $K(1s,2p)$ in eV. Vergelijk met $K(1s,2s) = \qty{0.398}{eV}$
en verklaar het verschil.
\end{exercise}

\begin{exercise}[$\star\star$]\label{exo:b3:many-electron-atoms:9}
Bepaal de termen van $d^2$ met de methode van dit hoofdstuk, vertrekkend
van de grootste $M_L$ (je hoeft niet de hele tabel uit te schrijven, maar
verantwoord elke term).
\end{exercise}

\begin{exercise}[$\star\star\star$]\label{exo:b3:many-electron-atoms:10}
Toon aan dat de ruimtelijke tripletfunctie $\frac{1}{\sqrt2}[1s(1)2s(2) -
2s(1)1s(2)]$ nul wordt wanneer $\mathbf r_1 = \mathbf r_2$ en dat de
singuletfunctie dat niet doet. Leg het verband met het teken van $E_+ - E_-$.
\end{exercise}

\begin{exercise}[$\star\star\star$]\label{exo:b3:many-electron-atoms:11}
% ledger: lev:Ni2.3F3, lev:Ni2.3F2
Voor \ce{Ni^2+} ($d^8$) liggen de niveaus ${}^3$F bij 0 ($J = 4$), 1360.7 ($J = 3$)
en \qty{2269.6}{cm^{-1}} ($J = 2$). Verklaar de volgorde, bereken de
verhouding van de intervallen en vergelijk met de voorspelling van Landé.
Geef $A$ uit het grootste interval.
\end{exercise}

\begin{exercise}[$\star\star\star$]\label{exo:b3:many-electron-atoms:12}
Toon aan dat de \omterm{def:b3:quantum-model-systems:hermitian}{verwachtingswaarde} van $\hat{\mathbf L}\cdot\hat{\mathbf S}$,
gesommeerd over alle $(2L+1)(2S+1)$ toestanden van een term, nul is, zodat
\omterm{def:b3:many-electron-atoms:spin-orbit}{spin-baankoppeling} het gewogen gemiddelde van de niveaus onveranderd
laat. Controleer dit voor de term ${}^3$P van koolstof.
\end{exercise}

\section{Probleem: twee heliums}

\begin{problem}[{Weekendprobleem --- de singulet- en tripletfamilie van
helium: determinanten, de termen van overgangsmetaalionen, de
selectieregels die de twee families scheiden, en de gemeten
uitwisselingsintegraal}]
\label{pb:b3:many-electron-atoms:1}
% ledger: lev:He.2s3S1, lev:He.2s1S0, lev:He.2p3P2, lev:He.2p3P1, lev:He.2p3P0, lev:He.2p1P1, lev:Ti2.3F3, lev:Ti2.3F4, lev:Ti2.3P0, lev:Ti2.3P1, lev:Ti2.3P2
Gegevens (NIST, in \unit{cm^{-1}} boven de grondtoestand van elk deeltje). Helium:
$1s2s$ \termsym{3}{S}{1} 159855.97, \termsym{1}{S}{0} 166277.44; $1s2p$
\termsym{3}{P}{2,1,0} 169086.76, 169086.84, 169087.83, \termsym{1}{P}{1}
171134.89. \ce{Ti^2+}: \termsym{3}{F}{2,3,4} 0, 184.9, 420.4; \termsym{3}{P}{0,1,2}
10538.4, 10603.6, 10721.2. $\qty{1}{eV} = \qty{8065.54}{cm^{-1}}$.

\medskip\noindent\textbf{Deel I --- Spinorbitalen en determinanten.}
\begin{enumerate}
  \item Geef de vier spinorbitalen die beschikbaar zijn voor de configuratie $1s2s$.
  \item Schrijf de vier determinanten op met één elektron in $1s$ en één in
        $2s$.
  \item Toon aan dat $|1s\alpha\,2s\alpha|$ het product is van een
        antisymmetrische ruimtelijke functie en de spinfunctie $\alpha(1)\alpha(2)$.
  \item Combineer $|1s\alpha\,2s\beta|$ en $|1s\beta\,2s\alpha|$ tot de
        component $M_S = 0$ van dezelfde ruimtelijke functie, en tot een vierde toestand.
  \item Welke ruimtelijke functie, symmetrisch of antisymmetrisch, hoort bij
        $S = 0$? En bij $S = 1$?
  \item Waarom heten deze toestanden singulet en triplet?
\end{enumerate}

\medskip\noindent\textbf{Deel II --- De termen van \ce{Ti^2+}.}
\begin{enumerate}[resume]
  \item Geef de configuratie van \ce{Ti^2+} en het aantal determinanten
        ervan.
  \item Bepaal de grondterm en het grondniveau met de regels van Hund;
        controleer aan de hand van de gegevens.
  \item De andere termen zijn ${}^3$P, ${}^1$G, ${}^1$D en ${}^1$S: controleer het aantal
        toestanden.
  \item Is de \omterm{def:b3:many-electron-atoms:spin-orbit}{fijnstructuur} van ${}^3$F normaal of omgekeerd? Toets de
        intervalregel van Landé.
  \item Bereken de gewogen gemiddelde energieën van de termen ${}^3$F en ${}^3$P.
  \item Hun afstand is $15B$, waarin $B$ een racahparameter van het ion is
        (\cref{ch:b3:complex-spectra-magnetism}). Bereken $B$.
\end{enumerate}

\medskip\noindent\textbf{Deel III --- Waarom twee families.}
\begin{enumerate}[resume]
  \item Welke selectieregel verbiedt een overgang tussen een singulet en een
        triplet?
  \item De laagste triplet, $1s2s$ \termsym{3}{S}{1}, kan niet uitzenden naar
        de grondtoestand $1s^2$ \termsym{1}{S}{0}. Geef twee redenen. Hoe
        heet zo'n toestand?
  \item Bereken de golflengte van de lijn $1s2p\ {}^3$P $\to 1s2s\ {}^3$S (gebruik
        het gemiddelde van de niveaus ${}^3$P).
  \item Bereken de golflengte van $1s2p\ {}^1$P $\to 1s2s\ {}^1$S.
  \item Bereken de golflengte van $1s2p\ {}^1$P $\to 1s^2\ {}^1$S. In welk
        gebied van het spectrum ligt ze?
  \item Verklaar waarom spectroscopisten in de negentiende eeuw, toen ze
        deze lijnen zagen, in twee verschillende gassen geloofden.
\end{enumerate}

\medskip\noindent\textbf{Deel IV --- De \omterm{def:b3:many-electron-atoms:exchange}{uitwisselingsintegraal}, gemeten.}
\begin{enumerate}[resume]
  \item Schrijf de eersteordeenergieën van de singulet en de triplet $1s2s$
        uitgedrukt in $E_{1s}$, $E_{2s}$, $J$ en $K$.
  \item Welke ligt lager, en waarom, in termen van de posities van de
        elektronen?
  \item Bereken de singulet-tripletafstand van $1s2s$ in \unit{cm^{-1}} en eV.
  \item Bereken $K(1s,2p)$ in eV uit de niveaus $1s2p$.
  \item Waarom is $K(1s,2s)$ groter dan $K(1s,2p)$?
  \item Wordt de eerste regel van Hund bevestigd door de aangeslagen
        toestanden van helium?
  \item Formuleer het resultaat: de \omterm{def:b3:many-electron-atoms:exchange}{uitwisselingsintegraal} $K(1s,2s)$ van
        helium, in eV.
\end{enumerate}
\end{problem}
